% ============================================================
% GUIDANCE SYSTEM AND PART A REUSE  (Requirement 4d)
% ============================================================
\section{Guidance System}
\label{sec:guidance}

\subsection{Subsystems Reused from Part A}
All four Part~A subsystems were carried into Part~B
(Table~\ref{tab:reuse}). They are described in the Part~A reports
\cite{partA_est, partA_group}; only the changes are given here.

\begin{table}[!t]
    \centering
    \caption{Part A Subsystems Reused in Part B}
    \label{tab:reuse}
    \footnotesize
    \begin{tabularx}{\linewidth}{@{}l c >{\raggedright\arraybackslash}X@{}}
        \toprule
        \textbf{Subsystem} & \textbf{Reused} & \textbf{Modification} \\
        \midrule
        State estimator   & Yes    & Same seven-state linear Kalman filter at $100$~Hz. Each pose sample is used in one update, on the step it arrives. $R$ and $Q$ recalibrated from the measured noise (Section~\ref{sec:estimator}). \\
        Waypoint guidance & Partly & Capture logic kept. Straight-line tracking replaced by following the planned path. \\
        Controller        & Yes    & Same structure. Speed limit halved to $1.0$~m/s; speed also scaled by the avoidance layer. \\
        Waypoint ordering & Yes    & Same ant colony optimiser. Cost is A* path length; five waypoints. \\
        \bottomrule
    \end{tabularx}
\end{table}

\subsection{State Estimator and Noise Characterisation}
\label{sec:estimator}
Every other subsystem acts on the estimate $\hat{\mathbf{x}}$, and the
obstacle positions are built from it (Eq.~\eqref{eq:detection}), so its
error sets the margins used throughout this report. The estimator is
the linear Kalman filter of Part~A \cite{partA_est}, run at $100$~Hz on
the state $\mathbf{x} = [p_{x}\ p_{y}\ \theta\ v\ \omega\ v_{x}\ v_{y}]^{\top}$.
Position integrates the Cartesian velocity and heading integrates the
turn rate, so the model is linear:
\begin{equation}
    \mathbf{x}_{k+1} = A\mathbf{x}_{k} + \mathbf{w}_{k}, \qquad
    \mathbf{z}_{k} = H\mathbf{x}_{k} + \mathbf{n}_{k},
    \label{eq:kf}
\end{equation}
where $A$ adds $\Delta t\,v_{x}$, $\Delta t\,v_{y}$ and $\Delta t\,\omega$
to $p_{x}$, $p_{y}$ and $\theta$, and $H$ selects the measured pose and
speed, $\mathbf{z} = [x\ y\ \theta\ v]^{\top}$. The velocity and
turn-rate states are random walks: a white disturbance of intensity $q$
entering a rate gives, for each (position, rate) pair,
\begin{equation}
\begin{aligned}
    Q_{\text{pair}} &= q\begin{bmatrix} \Delta t^{3}/3 & \Delta t^{2}/2 \\ \Delta t^{2}/2 & \Delta t \end{bmatrix}, \\
    R &= \operatorname{diag}(\sigma_{xy}^{2}, \sigma_{xy}^{2}, \sigma_{\theta}^{2}, \sigma_{v}^{2}).
\end{aligned}
\label{eq:kfnoise}
\end{equation}
The pose arrives at $20$~Hz. Each sample is used in one full update on
the step it arrives; on the steps between, only the speed is fused, so a
held sample is never counted twice. The covariance is updated in Joseph
form.

\textbf{Noise characterisation.} $R$ must describe the noise already
present in the simulator's signals. Part~A set it from the project
specification, variances of $0.1$~m$^{2}$ in position and $2$~deg$^{2}$
in heading, which matched the raw errors measured there, $0.32$~m and
$1.4^{\circ}$ \cite{partA_est}. That $R$ was carried into Part~B. The
Part~B model also outputs the vehicle pose before noise is added, so
the noise can be measured directly as the measured minus the noise-free
pose. Over the $25$ missions of Section~\ref{sec:results}
(Table~\ref{tab:noise}) it is zero-mean, uncorrelated between axes and
Gaussian in shape, with a standard deviation of $0.10$~m per axis and
$1.96^{\circ}$ in heading. The
Part~B simulator is therefore three times less noisy in position than
Part~A and somewhat noisier in heading, and the carried-over $R$ was
ten times too large in position. Consecutive samples are correlated ($0.5$) and samples
$0.1$~s apart are not; the filter treats every $20$~Hz sample as
independent, which slightly overstates the information it receives.

\begin{table}[!t]
    \centering
    \caption{Measured Sensor Noise, 25 Missions}
    \label{tab:noise}
    \footnotesize
    \begin{tabularx}{\columnwidth}{@{}X c c c@{}}
        \toprule
        & $x$ & $y$ & $\theta$ \\
        \midrule
        Mean                           & $-0.001$ m & $-0.001$ m & $-0.06^{\circ}$ \\
        Standard deviation             & 0.100 m & 0.097 m & $1.96^{\circ}$ \\
        Autocorrelation, 0.05 s / 0.10 s & 0.50 / 0.00 & 0.52 / 0.03 & 0.51 / 0.02 \\
        Kurtosis (Gaussian: 3)         & 3.22 & 2.97 & 3.07 \\
        \midrule
        \multicolumn{4}{@{}l}{Correlation between $x$ and $y$: $0.02$; \ $59\,624$ pose samples at $20$~Hz} \\
        \multicolumn{4}{@{}l}{Forward speed: mean $-0.008$~m/s, standard deviation $0.105$~m/s} \\
        \bottomrule
    \end{tabularx}
\end{table}

The speed noise was measured the same way, against the speed
differenced from the noise-free pose. Its standard deviation,
$0.105$~m/s, confirms the $\sigma_{v} = 0.1$~m/s carried over from
Part~A; its correlation in time cannot be resolved, because the
reference speed is itself derived from the $20$~Hz pose.
$R$ was set to the measured values, $\sigma_{xy} = 0.10$~m and
$\sigma_{\theta} = 2^{\circ}$, with $\sigma_{v} = 0.1$~m/s unchanged.
The Kalman gain depends on the ratio of $Q$ to $R$, so the velocity
intensity was lowered with it, from $0.5$ to $0.05$~m$^{2}$/s$^{3}$, with
$q_{\omega} = 5$~rad$^{2}$/s$^{3}$. Both settings were chosen by
replaying the filter offline on the logged measurements of missions
$1$ to $8$ over $60$ combinations of $\sigma_{xy}$, $q_{v_{x}, v_{y}}$
and $q_{\omega}$, scored against the noise-free pose. Accuracy is flat
near the optimum: the old and new settings both give $0.079$~m RMS
position error, against $0.078$~m for the best combination tested. What
the recalibration fixes is consistency. A consistent filter has a mean
normalised innovation squared (NIS) equal to the pose dimension, $3$;
the old $R$ gave $0.85$, a filter that overstates its own uncertainty,
and the new one gives $2.5$. Covariance calibration was the main
weakness identified in Part~A \cite{partA_est}. The filter reduces the planar position error
from $0.139$~m RMS in the raw measurement to $0.079$~m, and $99$\% of
its errors are below $0.18$~m. That value is what the $0.20$~m tracking
margin of Section~\ref{sec:planning} and the $0.25$~m obstacle margin of
Section~\ref{sec:avoidance} are sized against. In closed loop on missions
$1$ to $8$ (with the earlier $0.8$~m capture radius), the recalibrated filter also raised the worst true
clearance slightly, from $0.57$ to $0.59$~m to an obstacle and from
$0.63$ to $0.66$~m to a wall.

\subsection{Modifications for Part B}
\textbf{Path following.} In Part~A no wall stood between waypoints.
Here the straight line to a waypoint may pass through a wall, so
guidance follows the planned path, a list of points about $0.1$~m
apart. Guidance keeps an index $k$ into the list and steers at point
$\mathbf{p}_{k}$:
\begin{equation}
    \theta_{\text{cmd}} = \operatorname{atan2}\!\left(p_{k,y} - \hat{y},\ p_{k,x} - \hat{x}\right).
    \label{eq:guidance}
\end{equation}
The index advances whenever the vehicle is within $1.0$~m of
$\mathbf{p}_{k}$, so the target stays about $1.0$~m ahead on the path.

\textbf{Waypoint capture.} The index may not pass the point where the
current waypoint lies on the path, so the vehicle cannot skip a waypoint
by following the path past it. A waypoint is captured within
$r_{c} = 0.4$~m. The last one uses $0.10$~m; on its final approach the
vehicle steers directly at it, and on capture guidance stops the
simulation so that the mission ends on the waypoint.

\textbf{Path replacement.} When the re-planner publishes a new path,
guidance adopts it on the next step. If the leg being driven was
re-planned, the index restarts at the beginning of the new leg, which
is the vehicle's own position at the time of re-planning.

\textbf{Controller.} The turn rate is proportional to the heading error
$e_{\theta} = \theta_{\text{safe}} - \hat{\theta}$, and the speed to the
distance to the waypoint, reduced in turns and near obstacles:
\begin{equation}
    \omega = K_{\theta}\,e_{\theta}, \qquad
    v = K_{d}\,d_{\text{wp}}\cos e_{\theta}\; v_{\text{scale}}.
    \label{eq:controller}
\end{equation}
The turn rate is limited to $\omega_{\max} = 2\pi/5$~rad/s and
$K_{d}\,d_{\text{wp}}$ to $v_{\max} = 1.0$~m/s, and the factor
$\cos e_{\theta}$ is kept between $0.1$ and $1$. The gains are
$K_{\theta} = 1.0$~s$^{-1}$ and $K_{d} = 0.75$~s$^{-1}$. The speed limit
is the main trade-off between mission time and safety and is examined
with the other results in Section~\ref{sec:results}.

% ============================================================
% RESULTS
% ============================================================
\section{Results}
\label{sec:results}

The system was evaluated on $25$ random missions, generated from
seeds $1$ to $25$ so that they can be repeated. Mission time excludes
planning and re-planning, during which the simulation clock is stopped.
Ground truth is the vehicle pose before the simulator adds measurement
noise (Section~\ref{sec:estimator}). Clearance, cross-track error and
closest approach are measured from it. Clearance by the state estimate,
which is what the vehicle itself believed, is given for comparison.
Clearance is measured from the vehicle's centre to a wall face or an
obstacle's surface and must be at least $0.40$~m.

\begin{table}[!t]
    \centering
    \caption{Results over 25 Random Missions}
    \label{tab:results}
    \footnotesize
    \begin{tabularx}{\columnwidth}{@{}X c c c@{}}
        \toprule
        \textbf{Metric} & \textbf{Mean} & \textbf{Min} & \textbf{Max} \\
        \midrule
        Missions completed                    & \multicolumn{3}{c}{25 of 25} \\
        Mission time (s)                      & 122.9 & 78.2 & 175.3 \\
        Planned path length (m)               & 107.6 & 66.6 & 157.3 \\
        Min.\ obstacle clearance (m)          & 0.72 & 0.58 & 1.01 \\
        Min.\ wall clearance (m)              & 0.76 & 0.65 & 0.88 \\
        \quad by estimate: obstacle (m)       & 0.66 & 0.44 & 0.94 \\
        \quad by estimate: wall (m)           & 0.68 & 0.59 & 0.82 \\
        Position error of estimate, RMS (m)   & 0.079 & 0.077 & 0.081 \\
        Capture error, waypoints 1--4 (m)     & 0.39 & 0.38 & 0.40 \\
        Capture error, final waypoint (m)     & 0.094 & 0.062 & 0.100 \\
        Final stop, true distance (m)         & 0.14 & 0.05 & 0.22 \\
        Closest approach, waypoints 1--4 (m)  & 0.12 & 0.05 & 0.23 \\
        RMS cross-track error (m)             & 0.16 & 0.14 & 0.18 \\
        Max.\ cross-track error (m)           & 0.43 & 0.36 & 0.49 \\
        Re-plans per mission                  & 3.2 & 1 & 6 \\
        Re-plan computation (s)               & 0.20 & 0.12 & 0.65 \\
        Obstacle repulsion active (s)         & 0.36 & 0 & 1.61 \\
        Wall repulsion active (s)             & 9.1 & 2.8 & 17.2 \\
        \bottomrule
    \end{tabularx}
\end{table}

\textbf{Completion and time.} All $25$ missions captured every waypoint
in order (Table~\ref{tab:results}). The mean mission took $123$~s for a
$108$~m planned path. The vehicle drove $111$~m, $3.0$\% further, which
is the cost of the detours, of tracking error and of turning onto each
waypoint, at a mean speed of $0.90$~m/s against the $1.0$~m/s limit. By
its estimate the vehicle captured the last waypoint within $0.10$~m in
every mission. By the true pose it stopped $0.14$~m from it on average
and never more than $0.22$~m away; the difference is the estimate error
at the moment of capture. Figure~\ref{fig:mission} shows one mission.

\begin{figure}[!t]
    \centering
    \reportfig{fig_mission.png}
    \caption{Mission 9 as driven, by the true pose. Diamonds mark where
    the path was re-planned.}
    \label{fig:mission}
\end{figure}

\textbf{Clearance.} No mission breached the required clearance
(Fig.~\ref{fig:batch}). The closest true pass was $0.58$~m from an
obstacle, in mission~23, and $0.65$~m from a wall, in mission~14, a
margin of at least $0.18$~m in every mission. The vehicle's own
estimate never put it inside the limit either: its closest estimated
pass was $0.44$~m from an obstacle, in mission~9, where the true
clearance was $0.58$~m. The estimate can place the vehicle up to
$0.14$~m closer to an obstacle than it really is, and this is why the
margins of Sections~\ref{sec:planning} and~\ref{sec:avoidance} are
sized against the measured estimate error rather than against zero.

\begin{figure*}[!t]
    \centering
    \reportfig[0.9\textwidth]{fig_batch.png}
    \caption{Minimum true clearance (left) and mission time with number
    of re-plans (right) for each of the $25$ missions.}
    \label{fig:batch}
\end{figure*}

\textbf{Avoidance layers.} The path was re-planned between one and six
times in every mission. Each re-plan took $0.20$~s to compute on
average and $0.65$~s at most. The obstacle repulsion of the reactive
layer acted in $14$ of the $25$ missions, for $0.36$~s per mission on
average and never for more than $1.61$~s. The wall repulsion acted in
every mission, for $9.1$~s on average, mostly as a light correction in
doorways and narrow corridors. With the re-planner switched off,
six of missions $1$ to $8$ breached the obstacle clearance, the worst
passing $0.27$~m from an obstacle, with the reactive layer acting for
$4.9$~s per mission. At its backup tuning the reactive layer cannot keep the
clearance alone; re-planning is what provides it.

\textbf{Tracking and capture.} The cross-track error, measured from the
true pose against the path that was active at each instant, had an RMS
of $0.16$~m and a largest value of $0.49$~m (Fig.~\ref{fig:timeline}).
The peaks exceed the $0.20$~m added to the wall inflation; it is the
soft clearance, which keeps the path about $0.9$~m from walls wherever
there is room, that absorbs them. Waypoints $1$ to $4$ are captured on
entering a $0.4$~m capture radius, so their capture error equals the
radius by design. The path runs through each waypoint, so the vehicle
carries on past it after capture: its true path passed within $0.12$~m
of them on average and never more than $0.23$~m away.

\begin{figure}[!tb]
    \centering
    \reportfig{fig_timeline.png}
    \caption{True clearance and cross-track error in mission 9. Dotted
    lines mark the two re-plans.}
    \label{fig:timeline}
\end{figure}

\textbf{Speed and capture radius.} Table~\ref{tab:sweep} repeats
missions $1$ to $8$ with one parameter changed at a time. Raising the
speed limit to $1.5$~m/s cut the mission time by $23$\%, but one
mission passed $0.37$~m from an obstacle, and at $2.0$~m/s one mission
breached the clearance again. One cause is known: the paths are planned
for a turn radius of $1.0$~m, but at these speeds the tightest turn the
vehicle can hold grows to $1.2$~m and $1.6$~m, so it can no longer
follow the planned curves, and the RMS cross-track error rises from
$0.17$~m to $0.30$~m and $0.41$~m. A limit of $1.0$~m/s was therefore
the fastest of the three tested that kept the clearance, with no other
parameter re-tuned. Doubling the capture radius to $0.8$~m saved
$3.6$\% in time with no loss of clearance, but the vehicle then passed
the intermediate waypoints $0.35$~m away on average instead of
$0.13$~m. The $0.4$~m radius was chosen so that every waypoint is
visibly visited, at a small cost in time.

\begin{table}[!t]
    \centering
    \caption{Effect of Speed Limit (m/s) and Capture Radius (m), Missions 1--8, by the True Pose}
    \label{tab:sweep}
    \footnotesize
    \begin{tabularx}{\columnwidth}{@{}X c c c c@{}}
        \toprule
        & & $v_{\max}$ & $v_{\max}$ & $r_{c}$ \\
        & \textbf{Base} & \textbf{1.5} & \textbf{2.0} & \textbf{0.8} \\
        \midrule
        Mean mission time (s)          & 117.9 & 90.6 & 79.4 & 113.7 \\
        Worst wall clearance (m)       & 0.68 & 0.54 & 0.42 & 0.66 \\
        Worst obstacle clearance (m)   & 0.59 & 0.37 & 0.35 & 0.59 \\
        Missions below $0.40$~m        & 0 & 1 & 1 & 0 \\
        RMS cross-track error (m)      & 0.17 & 0.30 & 0.41 & 0.16 \\
        Closest approach (m)           & 0.13 & 0.10 & 0.07 & 0.35 \\
        \bottomrule
    \end{tabularx}
\end{table}
